\section{Formula targets independently of their original depth}
\label{sec:formula-balancing}

Formula balancing is classical, following Spira~\cite{spira1971}.
We include the elementary construction needed here to make the resulting
sample bound explicit. This section does not claim a new circuit-balancing
theorem.

\begin{lemma}[Polynomial-size balancing of a Boolean formula]
\label{lem:formula-balancing}
An AND/OR/NOT formula with $s$ gates and literal occurrences has an
equivalent fan-in-two circuit of size polynomial in $s+2$ and depth
$O(\log(s+2))$, regardless of the original formula depth.
\end{lemma}

\begin{proof}
Push negations to the leaves using De Morgan's laws, eliminate double
negations, and simplify constants. This does not increase the number
of literal occurrences. Binarize any larger-fan-in AND/OR gates; this
adds only linearly many gates and does not change the leaf count.
A constant formula is immediate. Otherwise
let $n$ denote its number of literal occurrences, counting repeated
occurrences separately.

For $n\ge2$ there is a subformula $g$ with between $n/3$ and $2n/3$
literal occurrences. To find one, descend from the root while a child
has more than $2n/3$ leaves. At the last such node both children have
at most $2n/3$ leaves, and at least one has at least $n/3$ leaves.
Regard this one occurrence of $g$ as a formal Boolean input. After
substituting either value $b\in\{0,1\}$ and simplifying constants,
the remaining formula $F_b$ has at most $n-|g|\le2n/3$ literal
occurrences. The Shannon identity is
\[
 F=(g\wedge F_1)\vee(\neg g\wedge F_0).
\]
It is valid even if variables appearing in $g$ also occur elsewhere
in the original formula.

Recursively balance $g,F_0,F_1$ and join the three resulting circuits
by the four indicated Boolean gates. Their depths increase by at most
three, and each recursive subproblem has at most $2n/3$ occurrences.
Hence the depth is $O(\log(n+1))$. There are at most three recursive
subproblems per call, so the size is polynomial; for example a bound
$O((n+1)^4)$ follows immediately by induction. Reuse the balanced
output of $g$ in its two appearances in the displayed identity.
This establishes the asserted circuit bounds.
\end{proof}

\begin{corollary}[A quasipolynomial bound for arbitrary formulas]
\label{cor:arbitrary-formula-learning}
Let $f$ have a Boolean formula with at most $s$ gates and literal
occurrences, and put $S=\max\{16,2+d+s\}$. There is a tanh realization
with unit margin, polynomial-in-$S$ width and physical weight bounds,
and depth $q$ satisfying
\[
 2\le q\le C\frac{\log S}{\log\log S}
\]
for a universal constant $C$. For each fixed $\zeta>0$ and every
learner depth $L\ge q$, the unrestricted posterior sample bound is
\[
 \exp\!\left[O_\zeta\!\left(\frac{\log^2S}{\log\log S}\right)\right]
 \bigl(L+\varepsilon^{-1}+\log(2/\delta)\bigr)^{O_\zeta(1)},
\]
simultaneously for all actual widths $N\ge C_L(f)$ and all
$0<\sigma^2\le N^{-(1+\zeta)}$.
This applies in particular to read-once formulas and does not depend
on the depth of the given formula.
\end{corollary}

\begin{proof}
Lemma~\ref{lem:formula-balancing} gives a circuit of polynomial size
and depth $O(\log S)$. Apply
Proposition~\ref{prop:boolean-layer-blocking} to that circuit.
The logarithm of its size parameter is $O(\log S)$ and its double
logarithm is comparable to $\log\log S$. Its blocked tanh depth is
therefore $O(\log S/\log\log S)$, with polynomial width and weights.
The quantified intrinsic-depth bound gives the displayed sample
estimate. All-width coverage is inherited from that bound, not from
identifying the target's actual minimum width with the construction.
\end{proof}

When a supplied formula is already shallower, its original blocked
representation can be used instead. The result above is still
quasipolynomial in formula size; it does not resolve polynomial sample
complexity for all polynomial-size formulas.
